\documentclass[10pt,a4paper]{amsart}
\usepackage{amsmath,amssymb,amsthm,graphicx}
\usepackage{verbatim}
\newenvironment{reference}{\comment}{\endcomment}
%\newenvironment{reference}{}{}
\newenvironment{slogan}{\comment}{\endcomment}
\newenvironment{history}{\comment}{\endcomment}
% For commutative diagrams we use Xy-pic
\usepackage[all]{xy}

% We use 2cell for 2-commutative diagrams.
\xyoption{2cell}
% Theorem environments.
%
\theoremstyle{plain}
\newtheorem{theorem}[subsection]{Theorem}
\newtheorem{proposition}[subsection]{Proposition}
\newtheorem{lemma}[subsection]{Lemma}

\theoremstyle{definition}
\newtheorem{definition}[subsection]{Definition}
\newtheorem{example}[subsection]{Example}
\newtheorem{exercise}[subsection]{Exercise}
\newtheorem{situation}[subsection]{Situation}

\theoremstyle{remark}
\newtheorem{remark}[subsection]{Remark}
\newtheorem{remarks}[subsection]{Remarks}

\numberwithin{equation}{subsection}

% Macros
%
\def\lim{\mathop{\mathrm{lim}}\nolimits}
\def\colim{\mathop{\mathrm{colim}}\nolimits}
\def\Spec{\mathop{\mathrm{Spec}}}
\def\Hom{\mathop{\mathrm{Hom}}\nolimits}
\def\Ext{\mathop{\mathrm{Ext}}\nolimits}
\def\SheafHom{\mathop{\mathcal{H}\!\mathit{om}}\nolimits}
\def\SheafExt{\mathop{\mathcal{E}\!\mathit{xt}}\nolimits}
\def\Sch{\mathit{Sch}}
\def\Mor{\mathop{\mathrm{Mor}}\nolimits}
\def\Ob{\mathop{\mathrm{Ob}}\nolimits}
\def\Sh{\mathop{\mathit{Sh}}\nolimits}
\def\NL{\mathop{N\!L}\nolimits}
\def\CH{\mathop{\mathrm{CH}}\nolimits}
\def\proetale{{pro\text{-}\acute{e}tale}}
\def\etale{{\acute{e}tale}}
\def\QCoh{\mathit{QCoh}}
\def\Ker{\mathop{\mathrm{Ker}}}
\def\Im{\mathop{\mathrm{Im}}}
\def\Coker{\mathop{\mathrm{Coker}}}
\def\Coim{\mathop{\mathrm{Coim}}}

% Boxtimes
%
\DeclareMathSymbol{\boxtimes}{\mathbin}{AMSa}{"02}

%
% Macros for moduli stacks/spaces
%
\def\QCohstack{\mathcal{QC}\!\mathit{oh}}
\def\Cohstack{\mathcal{C}\!\mathit{oh}}
\def\Spacesstack{\mathcal{S}\!\mathit{paces}}
\def\Quotfunctor{\mathrm{Quot}}
\def\Hilbfunctor{\mathrm{Hilb}}
\def\Curvesstack{\mathcal{C}\!\mathit{urves}}
\def\Polarizedstack{\mathcal{P}\!\mathit{olarized}}
\def\Complexesstack{\mathcal{C}\!\mathit{omplexes}}
% \Pic is the operator that assigns to X its picard group, usage \Pic(X)
% \Picardstack_{X/B} denotes the Picard stack of X over B
% \Picardfunctor_{X/B} denotes the Picard functor of X over B
\def\Pic{\mathop{\mathrm{Pic}}\nolimits}
\def\Picardstack{\mathcal{P}\!\mathit{ic}}
\def\Picardfunctor{\mathrm{Pic}}
\def\Deformationcategory{\mathcal{D}\!\mathit{ef}}
\usepackage{cleveref}
\crefname{originallemma}{original lemma}{original lemmas}
\Crefname{originallemma}{Original lemma}{Original lemmas}
\makeatletter
\newlabel{lemma-de-rham-cohomology-projective-space}{{0FMJ}{source}{}{sourcepointer-1}{}}
\newlabel{lemma-de-rham-cohomology-projective-space@cref}{{[originallemma][0][]0FMJ}{[1][0][]source}{}{}{}}
\makeatother
\crefname{originalproposition}{original proposition}{original propositions}
\Crefname{originalproposition}{Original proposition}{Original propositions}
\makeatletter
\newlabel{proposition-global-generation-on-fibres}{{0FMR}{source}{}{sourcepointer-2}{}}
\newlabel{proposition-global-generation-on-fibres@cref}{{[originalproposition][0][]0FMR}{[1][0][]source}{}{}{}}
\makeatother
\crefname{originalsection}{original section}{original sections}
\Crefname{originalsection}{Original section}{Original sections}
\makeatletter
\newlabel{section-projective-space}{{0FMG}{source}{}{sourcepointer-3}{}}
\newlabel{section-projective-space@cref}{{[originalsection][0][]0FMG}{[1][0][]source}{}{}{}}
\makeatother
\title{The de Rham cohomology of a projective bundle is free over the base, generated by powers of the first Chern class of the twisting sheaf}
\author{The Stacks Project Authors}
\date{stacks-project:derham.tex:local-snapshot-20261012; GNU FDL 1.2 or later}
\begin{document}
\maketitle
\noindent\textit{Editorial note.} All delivered mathematical text is the author's own, copied from the pinned Stacks Project chapter file derham.tex; no statement or proof step was rewritten, reordered or removed.\par\medskip
\noindent\textit{Editorial note.} Preamble provenance: the pinned chapter file opens with \textbackslash{}input\{preamble\} and carries no preamble text of its own, so the project's single Stacks-wide preamble staged at sources/stacks-derham/preamble.tex is excerpted verbatim as three spans, lines 8-12 (verbatim package and the project's own reference/slogan/history display environments), lines 14-18 (Xy-pic setup) and lines 154-217 (theorem declarations and mathematical macros). The wrapper class and the amsmath/amssymb/amsthm/graphicx base packages are supplied by the assembler. No macro was invented or redefined.\par\medskip
\noindent\textit{Editorial note.} Same-chapter pointers are located at their \textbackslash{}label lines inside the pinned chapter file; cross-chapter pointers are located in the project's own tag list staged in the snapshot (sources/stacks-chow/tags.txt), at the verbatim lines carrying their canonical tags.\par\medskip
\noindent\textit{Allowed standard prerequisites (editorial).}\begin{itemize}
\item Projective space and its basic properties. Same chapter derham.tex, Section \textbackslash{}ref\{section-projective-space\} (invoked at proof line 2268). It identifies the restriction of the projective bundle over a trivializing affine open with \$\textbackslash{}mathbf\{P\}\textasciicircum{}\{r-1\}\_A\$ and gives properness and smoothness of the projection.
\item The de Rham cohomology of projective space. Same chapter derham.tex, label lemma-de-rham-cohomology-projective-space (invoked at proof line 2272). It shows that the powers \$c\_1\textasciicircum{}\{dR\}(\textbackslash{}mathcal\{O\}\_P(1))\textasciicircum{}i\$, \$i = 0, \textbackslash{}ldots, r-1\$, freely generate the de Rham cohomology of each fibre over its residue field.
\item The criterion producing freeness over the base from a basis on every fibre. Same chapter derham.tex, label proposition-global-generation-on-fibres (invoked at proof line 2273). Applied to the proper smooth projection with the classes above, it turns the fibrewise basis statement into the asserted isomorphism.
\end{itemize}\medskip
\ifdefined\hypertarget\hypertarget{sourcepointer-1}{}\fi
\noindent\textit{Source reference.} Original lemma 0FMJ: sources/stacks-derham/pinned-original-source.tex, source lines 1802--1802. Navigation pointer for \textbackslash{}ref\{lemma-de-rham-cohomology-projective-space\} used inside the counted span (line 2272); the argument that consumes it is retained in the span.\par\medskip
\ifdefined\hypertarget\hypertarget{sourcepointer-2}{}\fi
\noindent\textit{Source reference.} Original proposition 0FMR: sources/stacks-derham/pinned-original-source.tex, source lines 2158--2158. Navigation pointer for \textbackslash{}ref\{proposition-global-generation-on-fibres\} used inside the counted span (line 2273); the argument that consumes it is retained in the span.\par\medskip
\ifdefined\hypertarget\hypertarget{sourcepointer-3}{}\fi
\noindent\textit{Source reference.} Original section 0FMG: sources/stacks-derham/pinned-original-source.tex, source lines 1576--1576. Navigation pointer for \textbackslash{}ref\{section-projective-space\} used inside the counted span (line 2268); the argument that consumes it is retained in the span.\par\medskip
\begin{proposition}
\label{proposition-projective-space-bundle-formula}
Let $X \to S$ be a morphism of schemes. Let $\mathcal{E}$ be a locally
free $\mathcal{O}_X$-module of constant rank $r$. Consider the morphism
$p : P = \mathbf{P}(\mathcal{E}) \to X$.
Then the map
$$
\bigoplus\nolimits_{i = 0, \ldots, r - 1} H^*_{dR}(X/S)
\longrightarrow
H^*_{dR}(P/S)
$$
given by the rule
$$
(a_0, \ldots, a_{r - 1}) \longmapsto
\sum\nolimits_{i = 0, \ldots, r - 1} c_1^{dR}(\mathcal{O}_P(1))^i \cup p^*(a_i)
$$
is an isomorphism.
\end{proposition}

\begin{proof}
Choose an affine open $\Spec(A) \subset X$ such that $\mathcal{E}$ restricts
to the trivial locally free module $\mathcal{O}_{\Spec(A)}^{\oplus r}$.
Then $P \times_X \Spec(A) = \mathbf{P}^{r - 1}_A$. Thus we see that
$p$ is proper and smooth, see Section \ref{section-projective-space}.
Moreover, the classes $c_1^{dR}(\mathcal{O}_P(1))^i$, $i = 0, 1, \ldots, r - 1$
restricted to a fibre $X_y = \mathbf{P}^{r - 1}_y$ freely generate the
de Rham cohomology $H^*_{dR}(X_y/y)$ over $\kappa(y)$, see
Lemma \ref{lemma-de-rham-cohomology-projective-space}. Thus we've verified the
conditions of Proposition \ref{proposition-global-generation-on-fibres}
and we win.
\end{proof}

\end{document}
